Bagian ini menjelaskan teori yang menjadi dasar pengembangan Judol Detector.
Konsep yang dibahas meliputi ekstensi browser, DOM, pencocokan string,
\textit{regular expression}, Weighted Levenshtein Distance, OCR, serta
mekanisme penandaan hasil deteksi pada halaman web.

\section{Ekstensi Browser Chromium}

Ekstensi browser adalah perangkat lunak tambahan yang berjalan di dalam
peramban untuk menambah atau memodifikasi perilaku halaman web. Pada Chromium,
ekstensi didefinisikan melalui file \texttt{manifest.json} yang memuat
metadata, permission, script, dan konfigurasi antarmuka \cite{chrome_extensions}.

Manifest V3 memisahkan beberapa komponen utama:

\begin{itemize}
    \item \textbf{Popup}, yaitu antarmuka kecil yang muncul ketika pengguna
    menekan ikon ekstensi.
    \item \textbf{Content script}, yaitu script yang disisipkan ke halaman web
    dan dapat membaca atau memodifikasi DOM halaman.
    \item \textbf{Permission}, yaitu izin yang diperlukan ekstensi, misalnya
    izin membaca tab aktif, menyimpan data lokal, atau mengeksekusi script.
\end{itemize}

Pada Judol Detector, popup digunakan untuk menampilkan statistik dan kontrol
pemindaian, sedangkan content script digunakan untuk mengambil teks halaman,
menjalankan highlight, dan memberi efek blur pada konten yang terdeteksi.

\section{Document Object Model}

\textit{Document Object Model} (DOM) adalah representasi struktur dokumen HTML
dalam bentuk pohon hierarkis \cite{mdn_dom}. Setiap elemen HTML, teks, dan
atribut dapat direpresentasikan sebagai node. Melalui DOM, script dapat membaca
isi halaman, mencari elemen tertentu, dan memodifikasi tampilan halaman secara
dinamis.

Dalam Judol Detector, DOM digunakan untuk dua keperluan utama. Pertama, content
script menelusuri DOM untuk menemukan node teks yang terlihat dan relevan untuk
dipindai. Kedua, setelah sebuah teks terdeteksi, node teks tersebut diganti
dengan elemen \texttt{span} khusus agar hasil deteksi dapat diberi highlight,
tooltip, dan blur.

Tidak semua node teks perlu dipindai. Teks yang berada pada elemen seperti
\texttt{script}, \texttt{style}, \texttt{noscript}, \texttt{textarea}, dan
\texttt{input} diabaikan karena tidak merepresentasikan konten halaman yang
umumnya dibaca pengguna. Pemfilteran ini membuat proses pemindaian lebih fokus
pada teks yang benar-benar tampil pada halaman.

\section{Pencocokan String}

Pencocokan string adalah proses mencari kemunculan suatu pola dalam teks.
Jika diberikan teks $T$ dengan panjang $n$ dan pola $P$ dengan panjang $m$,
tujuan pencocokan string adalah menemukan indeks pada $T$ di mana $P$ muncul.
Persoalan ini merupakan dasar dari banyak aplikasi pemrosesan teks, termasuk
deteksi kata kunci \cite{cormen2009introduction}.

Dalam konteks deteksi judol, pola yang dicari adalah daftar kata kunci seperti
\texttt{slot}, \texttt{gacor}, \texttt{rtp}, \texttt{jackpot}, dan frasa lain
yang berkaitan dengan promosi judi daring. Pencocokan dapat dilakukan secara
eksak maupun tidak eksak. Pencocokan eksak hanya menerima string yang sama
persis, sedangkan pencocokan tidak eksak menerima string yang mirip berdasarkan
ukuran tertentu.

\section{Regular Expression}

\textit{Regular expression} (RegEx) adalah notasi untuk mendeskripsikan pola
teks. RegEx memungkinkan pencarian tidak hanya berdasarkan kata tetap, tetapi
juga berdasarkan struktur karakter tertentu \cite{friedl2006regex}. Contohnya,
pola huruf yang diikuti angka dapat digunakan untuk mendeteksi nama situs atau
kode promosi yang ditulis seperti \texttt{SLOT88}, \texttt{ZEUS222}, atau
\texttt{MAXWIN234}.

Pada Judol Detector, RegEx digunakan untuk mendeteksi pola kata yang terdiri
atas minimal dua huruf dan diikuti dua hingga tiga angka. Pola ini menangkap
format umum yang sering muncul pada konten promosi judol, terutama ketika nama
brand atau istilah promosi digabungkan dengan angka.

Secara konseptual, pola yang digunakan dapat dijelaskan sebagai berikut:

\begin{table}[H]
\centering
\caption{Komponen pola RegEx}
\label{tab:regex-pattern}
\begin{tabular}{|p{4cm}|p{9cm}|}
\hline
\textbf{Komponen} & \textbf{Makna} \\
\hline
Batas kiri token & Memastikan pola tidak berada di tengah kata lain. \\
\hline
Huruf minimal dua karakter & Menangkap bagian nama atau istilah promosi. \\
\hline
Angka dua sampai tiga digit & Menangkap akhiran numerik seperti \texttt{88}, \texttt{777}, atau \texttt{222}. \\
\hline
Batas kanan token & Memastikan pola berakhir sebagai token utuh. \\
\hline
\end{tabular}
\end{table}

\section{Levenshtein Distance}

\textit{Levenshtein Distance} adalah ukuran jarak antara dua string berdasarkan
jumlah minimum operasi edit yang diperlukan untuk mengubah satu string menjadi
string lain \cite{levenshtein1966}. Operasi edit yang umum digunakan adalah:

\begin{enumerate}
    \item penyisipan karakter,
    \item penghapusan karakter,
    \item penggantian karakter.
\end{enumerate}

Jika jarak edit antara dua string kecil, maka kedua string dianggap mirip.
Sebagai contoh, kata \texttt{gacor} dan \texttt{g4cor} hanya berbeda pada satu
karakter jika angka \texttt{4} dianggap sebagai pengganti huruf \texttt{a}.

\subsection{Weighted Levenshtein Distance}

Weighted Levenshtein Distance adalah variasi Levenshtein Distance yang memberi
bobot berbeda pada operasi edit tertentu. Pada Judol Detector, penggantian
karakter yang mirip secara visual diberi biaya lebih kecil dibandingkan
penggantian karakter biasa. Pendekatan ini digunakan karena banyak konten judol
memodifikasi kata dengan karakter yang tampak mirip.

\begin{table}[H]
\centering
\caption{Contoh karakter yang dianggap mirip secara visual}
\label{tab:visual-substitution}
\begin{tabular}{|p{4cm}|p{8cm}|}
\hline
\textbf{Huruf} & \textbf{Variasi visual} \\
\hline
\texttt{o} & \texttt{0} \\
\hline
\texttt{a} & \texttt{4} \\
\hline
\texttt{i} & \texttt{1}, \texttt{l}, \texttt{|} \\
\hline
\texttt{e} & \texttt{3} \\
\hline
\texttt{s} & \texttt{5}, \texttt{\$} \\
\hline
\texttt{b} & \texttt{8} \\
\hline
\texttt{g} & \texttt{6}, \texttt{9} \\
\hline
\texttt{t} & \texttt{7} \\
\hline
\end{tabular}
\end{table}

Nilai kemiripan dihitung dari jarak edit berbobot. Jika nilai kemiripan
melebihi ambang batas tertentu, token pada halaman dianggap cocok dengan kata
kunci. Pada implementasi program, ambang batas default yang digunakan adalah
0,78. Nilai ini dipilih agar sistem cukup sensitif terhadap variasi penulisan,
tetapi tetap membatasi jumlah hasil positif palsu.

\section{Boyer-Moore}

Algoritma Boyer-Moore adalah salah satu algoritma pencocokan string yang
bekerja dengan membandingkan karakter dari kanan ke kiri. Inti dari algoritma
ini adalah menggunakan \textit{bad character heuristic} untuk menentukan seberapa
jauh pola digeser ketika terjadi ketidakcocokan.

\begin{quotation}
Inti dari \textit{bad character heuristic} cukup simpel. Ketika kita
membandingkan pola dengan teks dan menemukan karakter yang tidak cocok, karakter
di teks tersebut disebut \textit{bad character} . Begitu
terjadi ketidakcocokan, pola tidak digeser sedikit demi sedikit, melainkan
langsung ``dilompatkan'' sejauh mungkin. Ada dua kemungkinan pergeseran:
pertama, pola digeser hingga karakter buruk di teks bisa cocok dengan
kemunculan terakhirnya di dalam pola; kedua, jika karakter tersebut sama
sekali tidak ada di pola, maka pola langsung dilewati melewati karakter yang
tidak cocok itu.
\cite{geeksforgeeks_boyer_moore}
\end{quotation}

Dengan heuristic ini, algoritma dapat melompati sejumlah posisi yang tidak
perlu diperiksa, sehingga pada kasus terbaik kompleksitasnya mendekati $O(n/m)$.


\subsection{Algoritma Rabin-Karp}

Berbeda dengan algoritma pencocokan string yang membandingkan karakter satu per
satu secara langsung, algoritma Rabin-Karp menggunakan fungsi hash untuk
mempercepat pencarian. Ide utamanya adalah menghitung nilai hash dari pola, lalu
membandingkannya dengan hash setiap jendela teks yang berukuran sama dengan
pola. Jika hash cocok, barulah dilakukan pemeriksaan karakter demi karakter
untuk memastikan tidak terjadi \textit{hash collision}
\cite{geeksforgeeks_rabin_karp}.

Untuk menghitung hash, setiap karakter dianggap sebagai digit dalam sistem
bilangan berbasis $d$ (biasanya 256 untuk ASCII), lalu nilai tersebut diambil
modulo sebuah bilangan prima $q$. Modulo ini berfungsi untuk mencegah
\textit{overflow} dan mengurangi kemungkinan tabrakan hash. Agar tidak perlu
menghitung ulang dari nol setiap kali jendela bergeser, algoritma ini
menggunakan teknik \textit{rolling hash}: karakter paling kiri dikeluarkan,
sisa nilai dikalikan ulang dengan basis, lalu karakter baru di kanan
dimasukkan. Dengan cara ini, hash jendela berikutnya bisa diperoleh dalam waktu
konstan.

Pada kasus rata-rata, algoritma ini berjalan dengan kompleksitas $O(n + m)$
karena setiap pergeseran hanya melakukan satu kali operasi aritmetika. Namun,
pada kasus terburuk ketika banyak tabrakan hash terjadi, kompleksitasnya bisa
naik menjadi $O(n \times m)$ karena setiap tabrakan harus diverifikasi secara
manual karakter per karakter.

\section{Tokenisasi Teks}

Sebelum pencocokan fuzzy dilakukan, teks perlu dipecah menjadi kandidat token.
Token adalah potongan teks yang berisi huruf, angka, garis bawah, atau beberapa
simbol yang relevan. Tokenisasi memungkinkan program membandingkan kata kunci
dengan bagian teks yang berukuran sebanding.

Untuk kata kunci yang terdiri atas beberapa kata, seperti \texttt{slot gacor}
atau \texttt{bonus deposit}, program membentuk jendela token dengan jumlah
token yang sama dengan jumlah token pada kata kunci. Dengan cara ini, frasa
multi-kata tetap dapat dibandingkan menggunakan Weighted Levenshtein.

\section{Deduplication dan Cache}

Pipeline deteksi dapat menghasilkan lebih dari satu hasil pada rentang teks
yang sama, misalnya ketika satu token cocok oleh RegEx dan juga cocok oleh
Weighted Levenshtein. Oleh karena itu, program melakukan deduplikasi berdasarkan
rentang indeks awal dan akhir dari hasil deteksi. Hasil yang memiliki rentang
sama hanya disimpan sekali.

Program juga menggunakan cache untuk menyimpan hasil pemindaian teks. Cache
menghindari penghitungan ulang terhadap teks, daftar kata kunci, dan opsi
deteksi yang sama. Dengan demikian, pemindaian berulang pada halaman yang sama
dapat dilakukan lebih efisien.

\section{Optical Character Recognition}

\textit{Optical Character Recognition} (OCR) adalah proses mengenali teks dari
gambar. OCR berguna ketika teks tidak tersedia sebagai node DOM, melainkan
menjadi bagian dari banner, poster, atau gambar promosi. Pada Judol Detector,
OCR dilakukan menggunakan Tesseract.js, yaitu port JavaScript dari mesin OCR
Tesseract \cite{tesseractjs}.

Proses OCR pada program terdiri atas beberapa tahap:

\begin{enumerate}
    \item memilih gambar atau elemen visual yang terlihat pada viewport;
    \item mengambil teks konteks dari atribut gambar, nama file, dan container
    terdekat;
    \item menjalankan pipeline deteksi pada teks konteks;
    \item jika belum terdeteksi dan elemen berupa gambar yang valid, menjalankan
    OCR dengan Tesseract.js;
    \item menjalankan pipeline deteksi pada teks hasil OCR;
    \item memberi penanda dan blur pada gambar yang dinyatakan positif.
\end{enumerate}

Pendekatan ini menggabungkan OCR dengan analisis konteks. Analisis konteks
penting karena beberapa gambar promosi memiliki nama game atau kata kunci pada
caption, atribut \texttt{alt}, atau nama file, sedangkan OCR murni tidak selalu
mampu membaca teks pada gambar dengan akurat.

\section{Highlight, Tooltip, dan Blur}

Setelah hasil deteksi ditemukan, program perlu menyampaikan informasi tersebut
kepada pengguna. Judol Detector melakukan hal ini dengan memodifikasi DOM
halaman. Teks yang terdeteksi dibungkus dalam elemen \texttt{span} dengan class
khusus. Class tersebut diberi style agar teks terlihat sebagai highlight.

Setiap highlight menyimpan metadata deteksi pada atribut \texttt{data-*}, yaitu
kata kunci, algoritma, jumlah hasil, dan waktu eksekusi. Metadata ini digunakan
untuk menampilkan tooltip ketika pengguna mengarahkan kursor ke hasil deteksi.

Fitur blur digunakan untuk menyamarkan konten yang terdeteksi. Pada teks, blur
dapat diaktifkan atau dimatikan melalui popup. Pada gambar hasil OCR yang
positif, blur diberikan secara otomatis agar konten visual yang dicurigai tidak
langsung terlihat jelas.

\section{Kompleksitas Algoritma}

Analisis kompleksitas digunakan untuk memperkirakan biaya komputasi pipeline
deteksi. Jika $n$ adalah panjang teks, $k$ adalah jumlah kata kunci, dan $m$
adalah panjang rata-rata kata kunci, maka:

\begin{itemize}
    \item RegEx berjalan terhadap teks dengan biaya yang bergantung pada engine
    RegExp dan pola yang digunakan, tetapi pada pola sederhana umumnya linear
    terhadap panjang teks.
    \item Weighted Levenshtein membandingkan kandidat token dengan kata kunci.
    Setiap perbandingan dua string membutuhkan $O(m^2)$ pada kasus umum karena
    menggunakan tabel dinamis dua dimensi.
    \item OCR memiliki biaya paling besar karena melibatkan pengenalan citra.
    Oleh karena itu, jumlah gambar yang diproses dibatasi dan hasil OCR
    disimpan dalam cache.
\end{itemize}

Dalam implementasi, performa dijaga dengan beberapa strategi: membatasi jumlah
gambar OCR, menyaring gambar berdasarkan ukuran dan visibilitas, menggunakan
cache hasil teks dan OCR, serta melakukan deduplikasi hasil deteksi agar output
tidak berulang.
